\chapter{Deskripsi Persoalan}

Pemuatan barang ke dalam kendaraan angkut merupakan salah satu permasalahan umum dalam industri logistik yang melibatkan berbagai variabel, seperti dimensi paket, berat, nilai barang, kerapuhan, keterbatasan ruang, dan kapasitas kendaraan. Penentuan paket yang dimuat dan penempatannya dalam peti kemas truk memerlukan pengaturan yang efisien agar tidak terjadi pelanggaran, seperti paket yang melayang atau menembus dinding peti kemas, barang pecah-belah yang tertindih paket lain, atau total berat muatan yang melebihi kapasitas truk. Apabila dilakukan secara manual, proses ini sangat tidak efektif karena ruang solusi yang meningkat secara eksponensial seiring bertambahnya jumlah paket, dengan kombinasi pilihan paket, posisi koordinat, dan enam kemungkinan orientasi untuk setiap paket.

Tugas besar 1 mata kuliah IF3170 Inteligensi Artifisial bertujuan agar peserta kuliah memahami dan mampu mengimplementasikan serta mengevaluasi berbagai algoritma \textit{local search} untuk menyelesaikan permasalahan pengisian paket pada truk peti kemas Mejiro Express. Persoalan ini merupakan variasi dari \textit{3D bin packing} dan \textit{knapsack problem}, dengan tujuan memaksimalkan total nilai paket yang berada dalam peti kemas. Tidak seluruh paket harus dimuat, sehingga paket yang tidak terpilih dapat ditempatkan di luar truk. Dalam konteks inteligensi artifisial, proses ini dapat dimodelkan sebagai sebuah \textit{problem-solving agent} yang bertindak secara rasional dengan mempertimbangkan keterbatasan ruang dan kapasitas berat. Agen diharapkan dapat memilik aksi terbaik berdasarkan \textit{state} saat ini untuk mencapai kondisi optimal.

Untuk mencapai tujuan tersebut, digunakan pendekatan \textit{local search}. \textit{Local search} cocok untuk masalah yang tidak membutuhkan pencarian jalur menuju solusi dan hanya berorientasi pada pencarian \textit{final state} terbaik. Pencarian dimulai dari \textit{initial state} yang dibangkitkan secara acak, kemudian dieksplorasi terhadap \textit{state} tetangga untuk memperoleh nilai \textit{objective function} yang lebih baik. Perpindahan antar-\textit{state} dibatasi pada tiga jenis langkah, yaitu menukar dua paket, memindahkan sebuah paket ke koordinat yang berbeda, dan memutar paket 90 derajat terhadap sumbu x, y, atau z. Setiap \textit{state} harus memenuhi batasan yang ditetapkan, yaitu paket tidak boleh melayang atau menembus dinding peti kemas, barang pecah-belah tidak boleh menjadi penyangga paket lain, dan total berat muatan tidak boleh melebihi kapasitas truk.

Dalam tugas besar 1, dilakukan implementasi beberapa varian algoritma \textit{local search}, yaitu \textit{Hill-climbing, Simulated Annealing}, dan \textit{Genetic Algorithm}. Setiap algoritma memiliki karakteristik dan strategi pencarian yang berbeda. \textit{Hill-climbing} berfokus pada perbaikan bertahap menuju solusi terbaik, \textit{Simulated-Annealing} mengombinasikan \textit{Hill-climbing} dengan probabilitas berpindah ke solusi yang lebih buruk untuk menghindari \textit{local optimum}, serta \textit{Genetic Algorithm} yang mempertahankan sekumpulan solusi dan memperbaruinya melalui proses seleksi, \textit{crossover}, dan mutasi. Hasil dari setiap algoritma kemudian dianalisis berdasarlan nilai \textit{objective function} akhir, durasi pencarian, dan konsistensi hasil dari beberapa kali percobaan.